\subsection{叶内体积与法向输运}

对随参数变化的超曲面，前面的 Lie 输运公式可以在适配坐标中逐项落实。

令叶片诱导度规为 $h_{ab}$，单位法向为 $n$，约定 $K_{ab}=h_a{}^\mu h_b{}^\nu\nabla_\mu n_\nu$，$K=h^{ab}K_{ab}$。它与嵌入章法向分量 $B_{ab}$ 的关系是 $K_{ab}=-\langle B_{ab},n\rangle$。

只要能够在局部用一族超曲面描述空间，场的微分方程就会自然分成沿法向的变化与沿超曲面的变化。此时需要进一步考察法向、诱导度规和体积元之间的关系。

定义第二基本形式 $K_{ab}=\frac12(\mathcal L_n h)_{ab}$，并约定 $K=h^{ab}K_{ab}$。直接对 \cref{eq:ghs-foliation-adm} 求 Lie 导数可得
\begin{equation}\label{eq:ghs-extrinsic-curvature}
 K_{ab}=\frac1{2N}\bigl(\partial_t h_{ab}-D_aN_b-D_bN_a\bigr),
\end{equation}
其中 $D$ 是 $h$ 的 Levi--Civita 联络。这既是几何恒等式，又说明“叶的形变”不仅是坐标面积变化，还包含切向重标记。

\begin{proposition}[法向体积变化]
设 $X=X^n n+X^a\partial_a$，则
\[
 \nabla_\mu X^\mu=\nabla_nX^n+K X^n+D_aX^a
 - X^\mu_{\parallel}\,n^\nu\nabla_\nu n_\mu.
\]
特别地，$\nabla_\mu n^\mu=K$，故在切向位移为零、法向进动系数为一的 规范下 $\partial_t\sqrt h=K\sqrt h$。
\end{proposition}
\begin{proof}
由 $n_\mu n^\mu=1$ 得 $n^\mu\nabla_\nu n_\mu=0$。将单位切向投影 $h^\mu{}_{\nu}=\delta^\mu{}_\nu-n^\mu n_\nu$ 插入 $\nabla_\mu n^\mu$，得到 $h^{\mu\nu}\nabla_\mu n_\nu=K$。利用 $\partial_t\log\sqrt h=\frac12h^{ab}\partial_t h_{ab}$，再代入 \cref{eq:ghs-extrinsic-curvature}，可得一般式 $\partial_t\sqrt h=\sqrt h(NK+D_aN^a)$。令 $a_\mu=n^\nu\nabla_\nu n_\mu$。由于 $n_\mu X^\mu_\parallel=0$，沿 $n$ 求导给出
\[
 n^\mu n^\nu\nabla_\mu(X_\parallel)_\nu=-a_\nu X^\nu_\parallel,
 \qquad \nabla_\mu X^\mu_\parallel=D_aX^a-a_aX^a.
\]
而 $\nabla_\mu(X^nn^\mu)=\nabla_nX^n+KX^n$。把两式相加就得到散度分解；此处 $D_aX^a$ 是叶内散度，与环境全散度不同。
\end{proof}

对于切向位移为零 的具体情形 $g=\dd t^2+h_{ab}(t,y)\dd y^a\dd y^b$，直接从
$\Delta_g\varphi=(\sqrt g)^{-1}\partial_\mu(\sqrt g\,g^{\mu\nu}\partial_\nu\varphi)$ 算出
\begin{equation}\label{eq:ghs-laplace-evolving-slices}
 \Delta_g\varphi=\partial_t^2\varphi+K\partial_t\varphi+\Delta_{h(t)}\varphi.
\end{equation}
第二项 $K\partial_t\varphi$ 正是叶的体积膨胀对波动传播的影响。在球面上，它退化为熟悉的 $(d-1)r^{-1}\partial_r$；但该式在一般弯曲叶上依然成立。
